\begin{abstract}
Curvilinear structure segmentation is a challenging task: the targets are often
thin, low-contrast, and topologically fragile, while annotations are scarce and
test domains vary substantially.
Prompt-free adaptation of a frozen SAM with lightweight adapters has emerged as
a practical few-shot approach, but recent decoders often produce multiple
candidate masks and rely on test-time gating or fusion to select among them.
This input-dependent selection can be difficult to calibrate under domain shift,
can increase inference complexity, and may degrade structural connectivity.
For this family of prompt-free multi-candidate decoders, these issues suggest
that deployment can be a more important bottleneck than generating more
candidates.
We therefore propose Deterministic Candidate Decoding (DCD), which replaces
per-image selection with a validation-selected fixed deployment rule.
DCD trains four candidate heads with dense supervision and complementary
topology-aware objectives, but deploys a fixed rank (slot 3 in our benchmark)
at inference.
This reduces multi-candidate decoding to a deterministic single-output
deployment rule while retaining the benefits of candidate specialization.
Trained with only 18 images from six Seen datasets, DCD is evaluated on 12
datasets (8 Unseen and 4 Novel) spanning cracks, retinal and coronary vessels,
roads, wires, tires, and masonry mortar.
Against ten representative CNN-, transformer-, and SAM-based baselines, DCD
achieves the best macro-averaged performance among the compared methods on
this 12-dataset few-shot benchmark.
Averaged over four training seeds, it obtains 58.76 Dice, 45.17 IoU,
66.02 clDice, and 63.13 HD95, ranking first on all four macro metrics among
all compared methods.
Although the improvements are not uniform across every dataset, these results
expose the overlap--HD95 trade-off on cases such as DeepCrack and XCAD rather
than uniform superiority.
\end{abstract}
